% Oliver meeting, 2026-10-01 13:00 -- light progress update, adapted from
% slides_1005.tex's structure/theme (same content style, different
% audience: Oliver specifically, so this leads with the n=2->n=5 ANSYS/NEM
% work and the citation-attribution finding that concerns his own code).
% CONSTRAINT: keep this light -- aim for well under slides_1005.tex's
% 20-page cap; trim rather than pad.
\documentclass[aspectratio=169,11pt]{beamer}
\usetheme{Madrid}
\usecolortheme{seahorse}
\usepackage{booktabs}
\usepackage{amsmath,amssymb}
\usepackage{bm}
\usepackage{xcolor}
\usepackage{graphicx}

\definecolor{c3}{RGB}{30,80,160}
\definecolor{todo}{RGB}{200,30,30}
\setbeamercolor{frametitle}{bg=c3, fg=white}
\newcommand{\TODO}[1]{\textcolor{todo}{\textbf{[TODO: #1]}}}

\title{Master's Thesis --- Update for Oliver}
\subtitle{ANSYS/NEM ecology work since early September}
\author{Keisuke Nishioka\\[4pt]
\small Institute of Continuum Mechanics (IKM), Leibniz Universit\"at Hannover\\
\small Keio University}
\date{1 October 2026}

\begin{document}
\maketitle

% --------------------------------------------------------
\begin{frame}{Where things stand}
\begin{itemize}
  \item Admission granted 18 Aug 2026; submission planned early--mid
        December 2026
  \item FEM coupling route decided: \textbf{USERMAT} (ANSYS/Abaqus verified
        bit-identical, 0 ULP over 8017 states) with an optional embedded
        Python hook at the Gauss point
  \item This update: what changed on the \textbf{ANSYS/NEM side} (your
        pipeline) and the \textbf{USERMAT/ecology-bridge side} (this
        repo's own) since we last synced
\end{itemize}
\end{frame}

% --------------------------------------------------------
\begin{frame}{Headline finding: a citation mix-up in your interaction term}
\small
Checked \texttt{Interaction12}/\texttt{Interaction21} in your
\texttt{Ussfin} against the actual published equations directly (both
Klempt papers, read term-by-term) -- an earlier note of mine said this
term \emph{was} the 2026 paper's ``novel interaction scheme.'' That
was wrong.
\vspace{4pt}
\begin{center}
\begin{tabular}{@{}p{5.3cm}p{5.3cm}@{}}
\toprule
\textbf{Klempt et al.\ 2026 (AAM)}, Eq.\,10/16--18 &
\textbf{Your \texttt{Ussfin} term}\\
\midrule
Additive/bilinear: $I_{a,i}=\sum_j a_{ij}\bar\varphi_j$ added into a
dissipation residual &
Multiplicative: shifts a Monod-rate coefficient\\
$A$ symmetric, diagonal self-terms included &
No diagonal term, no symmetry constraint\\
Single scalar nutrient, linear &
Two independent Monod-saturated nutrient fields\\
\bottomrule
\end{tabular}
\end{center}
\vspace{4pt}
Your growth structure (Monod $\times$ Laplacian interface $\times$
chemotaxis) instead matches the \emph{other} paper, Klempt 2024 (BMMB) --
worth confirming which one you intended it to be.
\end{frame}

% --------------------------------------------------------
\begin{frame}{n=2 $\to$ n=5: generalized your interaction term, found a real limit}
\begin{itemize}
  \item Extended \texttt{Interaction12/21} to all 20 directed pairs among
        5 species (\texttt{usercm.inc}, \texttt{USolBeg}, \texttt{Ussfin})
        -- mechanical mirror of your own pattern, no formula change
  \item Ran on real ANSYS (\texttt{-smp -np 1} after any \texttt{Ussfin}
        recompile -- found that's required, DMP hangs silently)
  \item \textbf{Found a genuine critical coupling threshold}: species
        3$\leftrightarrow$4, your baseline constants
        (\texttt{MaxGrowth=100/dt=0.1/Penalty1=5}) --
        \textbf{0.0005 stable, 0.0007 diverges} (residual force overflow)
  \item Re-ran both boundary cases again 2026-09-08 to confirm it's real,
        not a fluke -- reproduced exactly
\end{itemize}
\vspace{4pt}
{\scriptsize Root cause: species 3's own uncoupled trajectory is already
close to unstable under your baseline constants -- any added coupling term
tips it over by the last substep. Full writeup:
\texttt{ansys\_usermat/apdl/N3\_GROWTH\_TRIAL.md}.}
\end{frame}

% --------------------------------------------------------
\begin{frame}{Separately: real CLSM data now drives this repo's own USERMAT}
\small Not your pipeline -- this repo's own ecology bridge
(\texttt{usermat\_biofilm.f}'s \texttt{kUseEcology}, the paper-faithful
$I_a=A\bar\varphi$ model). Replaced the placeholder seed with real Day-1
CLSM composition, step by step, each verified on real ANSYS against an
independent Python reference:
\vspace{4pt}
\begin{enumerate}
  \item Single fully-constrained element -- closed-form check, exact match
  \item Real unconstrained curved (cylinder) geometry -- exact match
  \item 4-region mesh, one real clinical condition per region
        (\textbf{CS/CH/DS/DH}, all four) -- exact match
\end{enumerate}
\vspace{4pt}
\texttt{psi} (viability) still a placeholder -- the workbook's ``living
cells'' ratio is a different, unbounded quantity from the model's
$\psi_i\in[0,1]$, confirmed by values $>1$ on this exact data; not
substituted rather than silently misused.
\end{frame}

% --------------------------------------------------------
\begin{frame}{UserElement vs.\ UserMat: still open}
\small The assignment says \emph{UserElement/UserMat}; this repo has a
UserMat only -- it can't transport the ecology field spatially by itself.
\vspace{6pt}
\begin{center}
\colorbox{c3!12}{\parbox{0.85\linewidth}{\centering\small
\textbf{Working assumption, pending you/Meisam:} extra nodal DOFs inside a
UserElement -- not precomputed externally, as this repo does today}}
\end{center}
\vspace{6pt}
{\scriptsize Our verified UMAT carries over either way
(\texttt{KEYANSMAT=1} $\to$ \texttt{ElemGetMat}, confirmed from ANSYS's
own example) -- the genuinely new work either way is the field-DOF
residual for ecology transport, no template in ANSYS's own examples.}
\end{frame}

% --------------------------------------------------------
\begin{frame}{Questions for you}
\begin{enumerate}
  \item Was \texttt{Interaction12/21} meant as the 2026 paper's scheme, or
        an independent addition on top of the 2024 growth structure? Which
        is the intended reference?
  \item Is a re-tuned parameter set (\texttt{MaxGrowth}/\texttt{dt}/
        \texttt{Penalty1}) worth exploring together, now that a concrete
        coupling-strength ceiling exists for the baseline one?
  \item UserElement vs.\ UserMat -- does the working assumption above
        still look right from your side?
\end{enumerate}
\end{frame}

\end{document}
